\xiaosi

自动驾驶规划系统负责根据感知信息生成安全、高效的驾驶轨迹，是自动驾驶技术栈中的核心环节。近年来，该领域经历了从传统模块化流水线向端到端学习的范式转变。UniAD\cite{hu2023uniad}、VAD\cite{jiang2023vad}、DriveTransformer\cite{drivetransformer2025}等端到端Transformer模型在nuScenes\cite{caesar2020nuscenes}等基准上取得了领先性能，推动了感知-预测-规划的联合优化。

然而，当前自动驾驶规划仍面临三大核心挑战。其一，\textbf{计算效率}问题：端到端大模型参数量巨大，难以在车端边缘设备上实时部署，而L2级自动驾驶系统要求端到端延迟在20--100ms之内。其二，\textbf{可解释性与安全合规}问题：欧盟《人工智能法案》将于2026年8月全面执行\cite{euaiact2025}，自动驾驶系统被归类为高风险AI系统，要求具备可追溯性和可解释性；当前端到端黑箱模型无法满足这一合规要求\cite{zablocki2022xad}。其三，\textbf{长尾场景泛化}问题：模型在分布外罕见场景下表现不稳定，难以保证安全性。

与此同时，生物启发的液态神经网络展示了独特潜力。Lechner、Hasani等人\cite{lechner2020ncp}提出的神经电路策略（Neural Circuit Policies, NCP）受秀丽隐杆线虫神经连接组\cite{white1986celegans}启发，仅用19个神经元和253个突触即可完成自动驾驶车道保持任务，参数量比传统LSTM网络少数个数量级，且注意力自动聚焦道路关键区域，具备内在可解释性。后续工作包括闭式连续时间网络CfC\cite{hasani2022cfc}（将推理速度提升5--8倍）和Liquid-S4\cite{hasani2023liquids4}（实现并行训练），最终演化为Liquid AI公司的LFM系列大模型\cite{lfm2tech2025}。

然而，NCP的四层布线拓扑（感觉层$\to$中间层$\to$命令层$\to$运动层）目前完全依赖手工设计，缺乏系统化的优化方法。此外，NCP在强化学习场景下的适用性尚未得到充分验证——其ODE动力学约束（权重非负、稀疏连接、固定极性）是否与Q-value估计兼容，是一个未被探索的问题。

基于上述背景，本文提出将NCP作为自动驾驶规划的轻量可解释Planning Head，设计进化搜索算法自动优化其布线拓扑，并在Highway-env\cite{leurent2018highway}仿真平台上进行系统验证。本研究旨在回答以下问题：（1）NCP能否以极少参数量达到甚至超越传统MLP/LSTM/GRU的决策性能？（2）能否通过自动化搜索发现比手工设计更优的NCP布线拓扑？（3）NCP的结构化稀疏连接是否带来可量化的性能和安全性优势？
